% ==========================================================================
%  notation.tex — every notational convention in the book lives here.
%  SHARED FILE: edit only in the "hub" chat (see CLAUDE.md).
%
%  Rule: chapters never type raw notation for anything with a macro below.
%  Changing a convention = editing one line here.
%  Keep the notation appendix (appendices/notation.tex) in sync.
% ==========================================================================

% ---- Number sets -------------------------------------------------------------
\newcommand{\N}{\mathbb{N}}          % natural numbers {1, 2, ...}
\newcommand{\Z}{\mathbb{Z}}
\newcommand{\Q}{\mathbb{Q}}
\newcommand{\R}{\mathbb{R}}
\newcommand{\C}{\mathbb{C}}

% ---- Delimiters (starred versions auto-size: \abs*{...}) ---------------------
\DeclarePairedDelimiter{\abs}{\lvert}{\rvert}
\DeclarePairedDelimiter{\norm}{\lVert}{\rVert}
\DeclarePairedDelimiter{\ceil}{\lceil}{\rceil}
\DeclarePairedDelimiter{\floor}{\lfloor}{\rfloor}
\DeclarePairedDelimiter{\inner}{\langle}{\rangle}
\newcommand{\set}[1]{\left\{ #1 \right\}}

% ---- Linear algebra ---------------------------------------------------------
\newcommand{\vect}[1]{\bm{#1}}       % vectors: bold, column by default
\newcommand{\mat}[1]{\bm{#1}}        % matrices: bold upper case
\newcommand{\T}{^{\top}}             % transpose: A\T
\DeclareMathOperator{\tr}{tr}
\DeclareMathOperator{\rank}{rank}
\DeclareMathOperator{\diag}{diag}
\DeclareMathOperator{\spn}{span}

% ---- Calculus and analysis --------------------------------------------------
\newcommand{\dd}{\mathop{}\!\mathrm{d}}  % upright differential: \int f(x) \dd x
\DeclareMathOperator*{\argmin}{arg\,min}
\DeclareMathOperator*{\argmax}{arg\,max}

% ---- Probability ------------------------------------------------------------
\newcommand{\Prob}{\mathbb{P}}
\DeclareMathOperator{\E}{\mathbb{E}}
\DeclareMathOperator{\Var}{Var}
\DeclareMathOperator{\Cov}{Cov}
\DeclareMathOperator{\Corr}{Corr}
\newcommand{\given}{\mid}                % conditioning: \Prob(A \given B)
\newcommand{\ind}{\mathbf{1}}            % indicator: \ind_A or \ind\{X > 0\}
\newcommand{\indep}{\perp\!\!\!\perp}    % independence
\newcommand{\iid}{\overset{\text{iid}}{\sim}}
\newcommand{\eqd}{\overset{d}{=}}        % equal in distribution
\newcommand{\convp}{\xrightarrow{\;p\;}}
\newcommand{\convd}{\xrightarrow{\;d\;}}
\newcommand{\convas}{\xrightarrow{\text{a.s.}}}

% ---- Named distributions (upright names) -------------------------------------
\newcommand{\Normal}{\mathrm{N}}         % \Normal(\mu, \sigma^2) — variance, not sd
\newcommand{\Bern}{\mathrm{Bernoulli}}
\newcommand{\Bin}{\mathrm{Bin}}
\newcommand{\Pois}{\mathrm{Poisson}}
\newcommand{\Geom}{\mathrm{Geom}}
\newcommand{\Unif}{\mathrm{U}}         % U[a, b] or U(a, b)
\newcommand{\Expo}{\mathrm{Exp}}         % Exp(\lambda) — rate parametrisation
\newcommand{\GammaDist}{\mathrm{Gamma}}  % Gamma(\alpha, \beta) — shape, rate
\newcommand{\BetaDist}{\mathrm{Beta}}
\newcommand{\NegBin}{\mathrm{NegBin}}     % NegBin(k, p): trials until the k-th success
\newcommand{\Hyp}{\mathrm{Hyp}}           % Hyp(n, r, N): red balls in n draws, r red of N
\newcommand{\ChiSq}[1]{\chi^2_{#1}}       % chi-squared with #1 degrees of freedom

% ---- Statistics ---------------------------------------------------------------
\DeclareMathOperator{\bias}{bias}
\DeclareMathOperator{\MSE}{MSE}
\DeclareMathOperator{\se}{se}
